\scanpage{6}
% archive foliation: 5
Dire que $f:\mathcal X_0\longrightarrow\mathcal X_1$ est \deleted{asphérique}\inserted{équiv. faible}, \unsure{signifie} aussi que
\diagram{\begin{tikzcd}
\mathcal X_0 \arrow[r] \arrow[r,shift right=2ex] \arrow[dr] & \mathcal X\\
& \mathcal X_1 \arrow[u,"\unsure{\sim}"']
\end{tikzcd}}
(imm. ouverte) l'est. (\uline{NB} Les faisceaux loc. constants sur $\mathcal X$ sont ceux de la forme $(F_1,f^*(F_1),u\ \text{tautol.})$, avec $F_1$ loc. ct. sur $\unsure{\mathcal X_0|\mathcal X_1}$.)

2) Partons inversement d'un topos $\mathcal X$, et d'un sous-topos ouvert
\[
(6)'\qquad\mathcal X_0\subset\mathcal X,
\]
$\unsure{\Longleftrightarrow|\Longleftarrow}$ sous-objet de l'objet final $e_{\mathcal X}$ de $\mathcal F(\mathcal X)$, d'où topos résiduel $\mathcal X_1$. On peut considérer aussi que la donnée de $\mathcal X_0$ correspond à celle d'un morphisme
\[
\pi:\mathcal X\longrightarrow\operatorname{Top}(\Delta_1)
\]
$\mathcal X_0$ étant l'image inverse du sous-topos ouvert ``universel'' $\{0\}$ dans $\operatorname{Top}(\Delta_1)$, et $\mathcal X_1$ de \deleted{\unsure{de}} l'image inverse du sous topos fermé $\{1\}$.
